\chapter{Теория}
\label{ch:intro}

\section{Обратное дискретное преобразование Фурье (ОДПФ)}

$$\boxed{s[n] = \frac{1}{N}\sum_{k=0}^{N-1}x[k]*e^{\frac{i2\pi nk}{N}}}$$

\section{Периодичность ДПФ}

ДВПФ периодична с периодом $2\pi$, а ДПФ периодична с периодом N, т.е

$$x[k+N] = x[k]$$

\section{ДПФ единичного импульса}

$$\boxed{x[k] = \sum{n=0}{N-1}\delta(n-n_0)e^{-i\frac{2\pi nk}{N}} = e^{-i\frac{2\pi n_0k}{N}}}$$

Свойства:

\begin{enumerate}
    \item |x[k]| = 1, $0 \le k \le N-1$ (равномерный спектр во всей полосе частот)
    \item $\phi[k] = -\frac{2\pi n_0k}{N}$
\end{enumerate}

\section{ДПФ гармонического колебания}
Основное применение ДПФ - спектральный анализ сигнала, в котором сигнал раскладывается на компоненты.\\

Рассмотрим ДПФ для анализа спектра гармонического колебания с частотой f, которое дискретизированно с частотой $f_s$

$$x(t) = cos(2\pi ft + \phi_0)$$

После дискретизации получим

$$x[n] = cos(2\pi(\frac{f}{f_s})n+\phi_0) = cos(\Omega n+\phi_0), 0 \le n \le N-1$$

Пусть N=16, f=1kHz, fs = 8kHz, $\phi_0 = -\frac{\pi}{2}$ \\

$$\Omega = \frac{1kHz}{8kHz}*2\pi = \frac{\pi}{4}$$

$$X[n] = cos(\frac{\pi n}{4} - \frac{\pi}{2})$$

$$T_s = \frac{1}{8kHz} = 0.125*10^{-3}с$$

Непрерывный сигнал дискретизируется на интвервале $NT_s$ секунд. Его $\Omega = \frac{\pi}{4} = \frac{2\pi}{8} = 2\frac{\pi}{N}$ на
N=16 точек соответствует двум периодом колебания x(t). Если нормированная частота колебания является целым кратным $\frac{2\pi}{N}$,
отсчеты ДПФ соответствуют значениям ДВПФ на частотах $\frac{2\pi l}{N}$, а на остальных частотах ДВПФ равно 0. В ином случае отсчеты
ДВПФ имеют ненулевые значения для всех индексов k. \\

Для спектрального анализа колебания применяют весовые окна (Ханна, Хэмминга)

$$\dot{X}(n) = X(n)*W(n)$$

Весовые окна применяются для снижения эффекта растекания спектра.

\endinput